% Avsnitt 10.10, ur lectures/L10/appendix/intro_framing_exercises.md och
% lectures/L10/appendix/c_exercises.md.

\section{Övningar}
\label{c10:sec:exercises}

\begin{aside}
Övningarna är upplagda så här. Övning~\ref{c10:ex:byhand} och \ref{c10:ex:cpp} hör till
\secref{c10:sec:framing}: den första bygger ramar för hand, på papper, och den andra gör samma ram
till en fungerande C++-implementation. Övning~\ref{c10:ex:ascan} slår bron mellan det ramformatet och
CAN. Övning~\ref{c10:ex:stuffing} till \ref{c10:ex:frame} befäster \secref{c10:sec:can} till
\ref{c10:sec:bittiming}; de görs på papper och bygger upp till att sätta ihop en hel ram för hand.
Övning~\ref{c10:ex:package} skriver paketet \code{can_def} och arbetar sedan mot det vid
tangentbordet.
\end{aside}

% ------------------------------------------------------------------------------------------
\exerciseset{Ramar för hand och i C++}

\exercise{Att konstruera två ramar för hand}{Resonemang}
\label{c10:ex:byhand}
Utgå från ramstrukturen ur \secref{c10:sec:framing}:

\begin{textdrawing}
Fält   Storlek   Beskrivning
------ --------- -------------------------------------------
SOF    2 byte    0xA5F7
LEN    1 byte    Nyttolastens längd
TYPE   1 byte    Meddelandetyp
DST    1 byte    Destinationsadress
SRC    1 byte    Källadress
SEQ    2 byte    Sekvensnummer
DATA   N byte    Nyttolast
CHK    2 byte    Checksumma (summan av föregående byte)
\end{textdrawing}

Checksumman är summan av varje byte från SOF till och med den sista nyttolastbyten, som en
\code{std::uint16_t} (så den slår runt vid overflow). Fält bredare än en byte (SOF, SEQ, CHK) är
big-endian. Ramtyperna är \code{Ping = 0x00}, \code{Pong = 0x01}, \code{StatusReq = 0x02} och
\code{StatusResp = 0x03}.

En enhet på adress \code{0x17} skickar en statusförfrågan till en sensor på adress \code{0x25}, med
sekvensnummer \code{0x7F05}. Det ger \code{TYPE = 0x02}, \code{DST = 0x25}, \code{SRC = 0x17},
\code{SEQ = 0x7F05}, och en tom \code{DATA}.

\xpart{a} Rita motsvarande \code{StatusReq}-ram på papper, och fyll i varje fält: SOF, LEN, TYPE,
DST, SRC, SEQ och CHK. Räkna ut checksumman för hand.

\xpart{b} Sensorn svarar med ett 16-bitars sensorvärde, \code{0x3201}, buret i nyttolasten. Rita
motsvarande \code{StatusResp}-ram på papper. Reglerna: TYPE ändras till \code{StatusResp}, DST och
SRC byter plats, SEQ är oförändrat, och DATA är \code{0x3201} (2 byte). Räkna ut även den här ramens
checksumma för hand.

\exercise{Ramen i C++}{C++}
\label{c10:ex:cpp}
Implementera ramen i C++ som en struct \code{comm::Frame}. Din implementation valideras av en
befintlig uppsättning enhetstester under \file{lectures/L10/appendix/code/}.

\task{1. Undersök filstrukturen}
\begin{itemize}
  \item \file{test/frame_test.cpp} innehåller enhetstesterna som validerar din ramimplementation. Dem
    skriver du inte; du får dem att passera.
  \item \file{include/comm/def.hpp} innehåller ramfältens offset (indexpositioner), fältstorlekarna
    med mera. Uppräkningsklassen \code{comm::FrameType} ska implementeras här.
  \item \file{include/comm/frame.hpp} innehåller deklarationen av structen \code{comm::Frame}.
  \item \file{source/comm/frame.cpp} ska innehålla implementationsdetaljerna för structen
    \code{comm::Frame}, alltså dess metoddefinitioner.
\end{itemize}

\task{2. Skapa ramtyperna}
Implementera uppräkningsklassen \code{comm::FrameType} i \file{comm/def.hpp}, och ge var och en av
de fyra typerna det id som listas i Övning~\ref{c10:ex:byhand}. Lägg till en avslutande uppräknare
\code{Unknown}, ett steg bortom den sista riktiga typen, så att en typ utanför intervallet kan
avvisas vid deserialisering med en enda jämförelse.

\task{3. Skapa ramstructen}
Lägg till medlemsvariablerna för structen \code{comm::Frame} i \file{comm/frame.hpp}, där
\code{@todo} markerar stället: \code{std::uint8_t data[MaxDataLen]}, \code{std::uint16_t seq},
\code{std::uint8_t dst}, \code{std::uint8_t src}, \code{std::uint8_t len} och \code{FrameType type}.

Värdeinitiera var och en med klammerparenteser: de numeriska medlemmarna med \code{\{\}}, och
\code{type} med \code{\{FrameType::Unknown\}}, så att en default-konstruerad ram aldrig håller en typ
som råkar vara giltig.

Bredderna är ingen smaksak. Enhetstesterna jämför en medlem mot en literal av fältets egen bredd
genom en mall som härleder \textbf{en} typ ur båda argumenten, så ett \code{len} deklarerat som
\code{std::size_t} eller ett \code{seq} deklarerat som \code{std::uint32_t} slösar inte bara
utrymme; testerna slutar kompilera.

\code{serialize()} och \code{deserialize()} är \textbf{redan deklarerade} i den filen, dokumenterade
och korrekt kvalificerade (\code{const noexcept} respektive \code{noexcept}). Låt de deklarationerna
vara; du skriver deras definitioner i steg 4 och 5.

\task{4. Implementera serialiseringsmetoden}
\Secref{c10:sec:framing}:s \code{checksum16()}, \code{writeToBuf()} och \code{readFromBuf()} ingår
\textbf{inte} i den utdelade koden, och båda metoderna nedan behöver alla tre. Kopiera in dem i
\file{comm/frame.cpp}, i en anonym namnrymd ovanför metoddefinitionerna, så att de förblir privata
för den filen:

\begin{cppcode}
namespace
{
// checksum16(), writeToBuf() and readFromBuf() from the framing introduction go here.
} // namespace
\end{cppcode}

Två saker om den kopieringen, som kompilatorn båda berättar för dig först i efterhand:
\begin{itemize}
  \item \code{writeToBuf()} och \code{readFromBuf()} är begränsade med \code{std::is_unsigned}, och
    ingenting i inkluderingskedjan för \file{frame.cpp} drar in \code{<type_traits>}. Lägg till den
    inkluderingen högst upp i filen.
  \item Ha \secref{c10:sec:framing}:s varning i åtanke när du anropar dem: \code{writeToBuf()}
    skriver \code{sizeof(T)} byte härledda ur \emph{argumentet}, inte ur fältet, så casta till
    fältets egen bredd vid varje anropsställe.
\end{itemize}

Implementera sedan metoden \code{comm::Frame::serialize()} så att den skriver följande till den
givna bufferten: \code{SOF = 0xA5F7} (big-endian), LEN, TYPE, DST, SRC, SEQ (big-endian), DATA i
samma ordning som den ligger i ramens nyttolastbuffert, och CHK, beräknad över alla föregående fält,
byte för byte.

Returnera ramens totala längd om serialiseringen lyckades, och \code{0U} vid fel, vilket betyder en
\code{buf} som är null eller en \code{bufLen} som är mindre än ramen som ska skrivas. Båda har ett
eget test; ingen av dem är valfri.

\task{5. Implementera deserialiseringsmetoden}
Implementera i samma fil metoden \code{comm::Frame::deserialize()} så att den validerar den givna
datan, \textbf{i den här ordningen}, eftersom varje kontroll gör nästa säker att utföra:
\begin{itemize}
  \item \code{buf} är inte null, och \code{bufLen} är minst \code{MinFrameLen}. Ingenting har lästs
    ännu vid den här punkten, så det är det här som över huvud taget gör det tillåtet att läsa
    huvudet. Hoppa över det och läsningarna av SOF och LEN nedan springer utanför en för kort
    buffert innan någon annan kontroll får en chans.
  \item SOF är korrekt.
  \item LEN ligger inom gränserna, alltså högst \code{MaxDataLen}.
  \item Den givna bufferten rymmer \textbf{hela} ramen som huvudet utger sig för. Det här är en
    andra, striktare längdkontroll, och den första ersätter den inte: en buffert som kapats efter
    nyttolasten lämnar ändå \emph{något} där CHK borde ligga, och att läsa det som checksumman får en
    trunkerad ram att se hel ut.
  \item Ramtypen är giltig, alltså att den mottagna byten ligger under \code{Unknown}.
  \item Checksumman stämmer.
\end{itemize}

Lagra den givna datan i ramen, men först när allt av den är giltigt.

Två av de punkterna är kontroller som testerna straffar hårdast. En saknad null-kontroll får inget
test att fallera, den \textbf{kraschar hela körningen}, så du ser en segfault och inget testnamn. En
buffertlängdskontroll som stannar vid huvudet fäller \code{DeserializeRejectsShortBuffer} och
\code{DeserializeRejectsTruncatedPayload} och ingenting annat, vilket är ett betydligt vänligare sätt
att få det sagt.

\task{6. Validera din implementation}
Enhetstesterna ligger bakom ett makro, så ingenting körs förrän du slår på det. Lägg till
\code{-DL01} i \code{CXX_FLAGS} i makefilen, där en \code{Todo}-kommentar markerar stället:

\begin{makecode}
CXX_FLAGS := -Wall -Werror -std=c++17 -Iinclude -DL01
\end{makecode}

Bygg och kör dem sedan från en Linuxterminal:

\begin{shell}
cd lectures/L10/appendix/code
make
\end{shell}

Tills du gör det rapporterar programmet att testerna är avstängda och avslutas med lyckat resultat,
oavsett om din implementation fungerar eller inte.

\begin{aside}
\textbf{Om bygget avbryts direkt med det här, är det inte din kod:}

\begin{console}
make[1]: *** No rule to make target 'lib'.  Stop.
make: *** [Makefile:26: ../../../../libs/test/libqacademy_test.a] Error 2
\end{console}

Testramverket ligger som en submodul i \file{libs/test}, och en klon utan
\code{--recurse-submodules} lämnar katalogen tom. Att den ändå \emph{finns} är varför felet handlar
om ett saknat make-mål och inte om en saknad katalog: det är ramverkets makefil som inte är hämtad
än. Hämta den en gång, från repots rot, med \code{git submodule update --init}. Övningen behöver
också en C++-kompilator.
\end{aside}

Testerna kontrollerar båda metoderna mot ramarna du byggde för hand i Övning~\ref{c10:ex:byhand}:
PING-bytesen ur \secref{c10:sec:framing} och båda ramarna ur den övningen, byte för byte; en rundtur
genom serialisering och deserialisering som måste bevara varje fält; och avvisningsfallen (felaktig
SOF, felaktig checksumma, ogiltig typ, trunkerad buffert, nullpekare).

Varje test som passerar skrivs ut i grönt; ett som fallerar skrivs ut i rött och namnger uttrycket,
båda värdena och raden, så att du kan jämföra direkt mot din egen pappersuträkning. Räkna med att se
alla fallera först, och gå gröna en grupp i taget medan du arbetar dig igenom steg 2--5.

Två saker som är värda att veta i förväg:
\begin{itemize}
  \item Testerna läser \code{frame.type}, \code{frame.dst}, \code{frame.src}, \code{frame.seq},
    \code{frame.len} och \code{frame.data} direkt, så de medlemsnamnen måste stämma. Har du döpt dina
    till något annat står det dig fritt att ändra testerna så att de matchar.
  \item \code{DeserializeLeavesFrameUntouchedOnFailure} ger \code{deserialize()} en korrupt buffert
    och kontrollerar sedan att ramens fält är \emph{oförändrade}. En \code{deserialize()} som
    kopierar in fält i ramen medan den parsar dem passerar alla andra tester och fäller det här.
\end{itemize}

% ------------------------------------------------------------------------------------------
\exerciseset{Från ramformatet till CAN}

\exercise{Samma ram, fast som CAN}{Resonemang}
\label{c10:ex:ascan}
Ta den \code{StatusResp}-ram du byggde för hand i Övning~\ref{c10:ex:byhand}\,b), och skicka den som
en CAN-ram i stället. Dess fält, som påminnelse: \code{TYPE = StatusResp}, \code{DST = 0x17},
\code{SRC = 0x25}, \code{SEQ = 0x7F05}, nyttolast 2 byte med värdet \code{0x3201}.

Utgå från konventionen i \secref{c10:sec:broadcast}: identifieraren kombinerar en bas per typ med
id:t för den sensornod meddelandet gäller, alltså den nod som frågas i en förfrågan och den nod som
svarar i ett svar (nod 3 frågas på \code{0x303} och svarar på \code{0x403}):

\[ \mathit{ID} = \mathit{TYPE\_BASE} + \mathit{nodeId} \]

Typbaserna är \code{StatusReq = 0x300}, \code{StatusResp = 0x400}, \code{Ping = 0x500} och
\code{Pong = 0x600}.

\xpart{a} Skriv motsvarande CAN-ram på formen \code{ID | DLC | DATA}.

\xpart{b} Ange vilka fält i \secref{c10:sec:framing}:s ram som saknar motsvarighet i CAN-ramen. Säg
för varje saknat fält vad CAN gör i stället, eller förklara varför CAN inte behöver något
motsvarande fält.

\xpart{c} Anta att två noder börjar sända på CAN-bussen i exakt samma ögonblick. Förklara kort hur
bussen avgör vilken nod som fortsätter sända, och hur den förlorande noden upptäcker att den har
förlorat. Det här är arbitreringsmekanismen från \secref{c10:sec:arbitration}; återge den med egna
ord.

\xpart{d} Förklara varför det inte finns något separat checksummefält att fylla i när man
konstruerar en CAN-ram, till skillnad från \secref{c10:sec:framing}:s ramformat.

% ------------------------------------------------------------------------------------------
\exerciseset{Bitstoppning, arbitrering och tajming}

\exercise{Bitstoppning för hand}{Resonemang}
\label{c10:ex:stuffing}
\xpart{a} Tillämpa stoppbitsregeln på följande 10-bitarssekvens:

\begin{textdrawing}
0011111111
\end{textdrawing}

Skriv den kompletta sända bitströmmen, med alla ursprungliga bitar och varje instoppad stoppbit
tydligt markerad.

\xpart{b} En mottagare observerar följande sända ström på 9 bitar:

\begin{textdrawing}
101111101
\end{textdrawing}

Mottagaren vet att bitstoppning är i bruk. Avstoppa strömmen: identifiera eventuella stoppbitar, ta
bort dem, och återskapa den ursprungliga följden av riktiga databitar.

\xpart{c} En mottagare observerar sex dominanta (\code{0}) bitar i rad på ett ställe där en stoppbit
skulle ha brutit en följd av fem. Förklara vad mottagaren drar för slutsats, och vad den bör göra
härnäst.

\xpart{d} \Secref{c10:sec:stuffing} påpekar att räknaren för lika bitar i rad inte nollställs vid
fältgränser. Anta att ett fält slutar med tre lika bitar och att nästa börjar med två till av samma
värde. Använd regeln och säg var stoppbiten hamnar, räknat från fältgränsen, och förklara varför
inget av fälten på egen hand skulle ha orsakat den.

\exercise{Att avkoda en riktig bitström}{Resonemang}
\label{c10:ex:decode}
En mottagare observerar följande ström på 21 bitar:

\begin{textdrawing}
0 00000010101 0 00 0010 11
\end{textdrawing}

Strömmen är redan avstoppad, så att den här övningen håller fokus på fältens uppdelning i stället
för på stoppningen. Mellanrummen finns bara med för läsbarhetens skull; den riktiga CAN-bussen har
inga glapp mellan fälten.

\xpart{a} Dela upp strömmen i dess namngivna fält, med bredderna från \secref{c10:sec:frameformat}:
SOF, ID, RTR, IDE, r0, DLC. Bestäm sedan identifieraren i hexadecimal form, DLC-värdet, och vilket
fält som har börjat i de två bitar som blir över när alla namngivna fält är avräknade.

\xpart{b} Bestäm utifrån DLC:n från \textbf{a)} hur många databitar mottagaren förväntar sig totalt,
och hur många ytterligare databitar som återstår innan CRC-fältet börjar.

\xpart{c} Anta att DLC i stället hade varit \code{0000}. Vilket fält skulle då följa omedelbart efter
kontrollfältet?

\exercise{Arbitrering mellan tre noder}{Resonemang}
\label{c10:ex:arb3}
Tre noder försöker sända vid samma SOF: nod R med ID \code{0x155}, nod S med ID \code{0x154}, och
nod T med ID \code{0x1D5}.

\xpart{a} Arbeta igenom arbitreringen bit för bit, i samma stil som \secref{c10:sec:arbitration}:s
genomräknade exempel. Bestäm vilken nod som vinner, och vid vilken identifierarbit varje förlorande
nod faller ur.

\xpart{b} De två förlorande noderna faller ur vid olika bitpositioner. Förklara varför genom att ange
var varje förlorande nods identifierare först skiljer sig från vinnarens identifierare, och varför
just den första avvikande biten avgör när noden förlorar arbitreringen.

\xpart{c} En förlorande nod måste sluta driva bussen, och den intressanta frågan är \emph{när}.
(\Chapref{c11:ch} introducerar portparet \code{tx_bus}/\code{bus_en} som bokens kontroller slutar
driva med, och \chapref{c16:ch}--\ref{c17:ch} bygger logiken; här resonerar du på protokollnivå.)
Förklara, med en eller två meningar, vad varje förlorande nods \code{can_controller} måste göra på
allra första klockcykeln efter att den upptäckt avvikelsen, och varför den tajmingen spelar roll.

\note[Ledtråd]Tänk igenom vad ”ögonblicket” betyder i \secref{c10:sec:arbitration}:s genomräknade
arbitreringsexempel.

\exercise{Resonemang om sampelpunkten}{Resonemang}
\label{c10:ex:sample}
Den här boken samplar varje bit 70~\% in i bitperioden, inte vid 0~\% (allra först i perioden) eller
100~\% (allra sist).

\xpart{a} Förklara vad som skulle kunna gå fel om en mottagare samplade vid 0~\%. Utgå från
signalutbredning, varför en bitperiod behöver en minsta varaktighet, och om det sända värdet har
hunnit stabilisera sig.

\xpart{b} Förklara vad som skulle kunna gå fel om en mottagare samplade vid 100~\%, på den sista
möjliga klockpulsen. Gör förklaringen specifik för att det i bokens design gäller att varje nod har
sin egen lokala bittimer, att timern omsynkroniseras bara vid SOF, och att den aldrig
omsynkroniseras mitt i ramen (riktig CAN omsynkroniserar på flanker även mitt i ramen; bokens
kontroller gör inte det).

\xpart{c} Arbitrering kräver att varje konkurrerande nod läser bussen under \emph{samma} bitperiod
som den driver den. Förklara varför en sampelpunkt för nära periodens början är riskabel, varför en
sampelpunkt för nära periodens slut är riskabel, och varför faran växer ju längre in i ramen man
kommer.

\exercise{Att konstruera en CAN-ram för hand}{Resonemang}
\label{c10:ex:frame}
Övning~\ref{c10:ex:stuffing} till \ref{c10:ex:sample} arbetade med enskilda regler. Den här sätter
ihop en komplett ram utifrån ramdiagrammet i \secref{c10:sec:frameformat}.

\begin{aside}
\textbf{En anmärkning om CRC:n.} Du får CRC-15-värdet för varje ram nedan. Att räkna ut en 15-bitars
CRC över drygt 30 bitar för hand är mekaniskt snarare än lärorikt; i \chapref{c13:ch} arbetar du
igenom rekursionen ordentligt, över en 4-bitarssekvens, när motorn i sig är ämnet.
\end{aside}

\xpart{a} Ta CAN-ramen du tog fram i Övning~\ref{c10:ex:ascan}: \code{ID = 0x425}, \code{DLC = 2},
data \code{32 01}. Skriv ut varje fält som bitar, i sändningsordning, med bredderna från
\secref{c10:sec:frameformat}: SOF; ID (11 bitar, MSB först); RTR; IDE och r0; DLC; data (MSB först,
byte 0 först); CRC (given: \code{010101100100011}); CRC-avgränsare, ACK-plats, ACK-avgränsare och
EOF.

\xpart{b} Hur många bitar blir det totalt, före all stoppning?

\xpart{c} Tillämpa nu stoppbitsregeln från Övning~\ref{c10:ex:stuffing} över SOF till och med
CRC-fältet, och kom ihåg att följdräknaren inte nollställs vid fältgränser. Hur många stoppbitar
sätts in, och hur många bitar upptar ramen på ledningen?

\xpart{d} Vid bokens bithastighet på 1~Mbit/s, hur lång tid tar det att sända den ramen?

\xpart{e} Upprepa \textbf{a)} till \textbf{d)} för en ram \textbf{utan nyttolast}:
\code{ID = 0x100}, \code{DLC = 0}, CRC \code{011100000001010}. Notera vilka delar av ramen som är
strukturellt oförändrade trots det saknade datafältet, och var CRC-fältet nu börjar.

\xpart{f} Jämför dina två svar på \textbf{c)}. Den andra ramen bär ingen data alls, och ändå behöver
dess stoppade område \emph{fler} stoppbitar än den första. Förklara varför, utifrån vad
stoppbitsregeln faktiskt reagerar på.

% ------------------------------------------------------------------------------------------
\exerciseset{Att arbeta med paketet}

\exercise{\code{can_def} och dess härledningar}{VHDL}
\label{c10:ex:package}
Varje konstant i \code{can_def} går att spåra tillbaka till något det här kapitlet definierat. De
här frågorna skriver paketet, kontrollerar att härledningarna håller när indata ändras, och utökar
det med värden som \secref{c10:sec:frameformat} namngav men som paketet ännu inte håller.

\xpart{a} Skriv \file{controller/can_def.vhd}, med konstanterna och subtyperna som
\secref{c10:sec:candef} ställer upp. Var du med på passet har du den redan; var du inte det är det
här den kommer ifrån, och varje kapitel från \chapref{c11:ch} och framåt förutsätter att den finns.

Det här är den enda modulen i boken som delas ut i sin helhet i stället för som en specifikation, och
det av ett skäl värt att nämna: ett paket innehåller inget beteende att designa, så det finns
ingenting mellan värdena och koden för dig att lista ut. Det nya här är själva VHDL-koden, eftersom
en paketdeklaration är en konstruktion boken inte använt tidigare. Skriv av den för hand i stället
för att klistra in den. Namnen är de kommande åtta kapitlens vokabulär, och att skriva dem en gång är
billigare än att slå upp dem åtta gånger.

\xpart{b} Anta att DE0-CV:ns klocka vore 40~MHz i stället för 50~MHz, med bithastigheten fortfarande
1~Mbit/s. Räkna om \code{TICKS_PER_BIT} och \code{SAMPLE_TICK} (fortfarande 70~\%).

\xpart{c} Anta i stället att klockan ligger kvar på 50~MHz men att bussen körs på 125~kbit/s, en
verklig CAN-hastighet för ”low speed”. Räkna om båda konstanterna.

\xpart{d} \code{TICKS_PER_BIT} är \code{CLOCK_FREQ_HZ / BIT_RATE_HZ}, med VHDL:s heltalsdivision för
\code{natural}, som trunkerar. Ge ett par av klocka och bithastighet där divisionen \emph{inte} går
jämnt ut, och förklara med en mening varför det vore ett problem för den här designen, trots att
VHDL-koden fortfarande skulle kompilera och elaboreras utan invändningar.

\xpart{e} Lägg till en konstant \code{BIT_PERIOD_NS}, längden på en CAN-bitperiod i nanosekunder,
härledd ur \code{BIT_RATE_HZ} i stället för hårdkodad. Ange sedan sambandet mellan den,
\code{TICKS_PER_BIT} och 50 MHz-systemklockans periodtid på 20~ns, och bekräfta att de tre stämmer
överens.

\xpart{f} \Secref{c10:sec:frameformat} anger CRC-avgränsaren, ACK-fältet och EOF som 1, 2 respektive
7 bitar, men paketet håller ingen av dem: \code{can_controller} kommer att använda de talen som nakna
literaler. Lägg till \code{CRC_DELIM_BITS}, \code{ACK_BITS} och \code{EOF_BITS}, och säg vilken av de
två formerna en senare läsare är bäst betjänt av.

\xpart{g} Kontrollera att paketet faktiskt kompilerar. Installera GHDL om du inte redan gjort det
(bilaga~\ref{app:sim}), och kör sedan \code{make build-project} från roten av gruppens repo. Du bör
se:

\begin{console}
--> controller/can_def.vhd analyzes cleanly
\end{console}

följt av de två utdelade \file{bridge/}-filerna, och därefter varje testbänk rapporterad som
överhoppad, eftersom modulerna de kontrollerar inte finns ännu:

\begin{console}
0 testbench(es) run, 8 skipped.
\end{console}

Gör nu sönder något med flit, ett saknat semikolon eller en felstavad typ, kör om, och läs vad GHDL
säger om det. Ställ tillbaka det.

Det här är hela verifieringsberättelsen för ett paket: ett paket har inget beteende att simulera, så
”den analyseras” är allt som finns att kontrollera. \Chapref{c11:ch} introducerar verktyget
ordentligt, och därifrån får varje modul en testbänk också.
